\cotefile{29}
\scanpage{157}
\uline{Proposition 3.1.} (\unsure{Remontes|Remontées|Réunions}) $X$ étale sur $Y$ (par ex. fini).
$G$ groupe fini \deleted{\unsure{de|dé|dans}} \inserted{opérant par} $Y$-automorphismes sur $X$.

\deleted{Si $s\in$} \inserted{(i)} L'ens des pts de $X$ où $G$ opère
\inserted{\unsure{ou s'il y a lieu|où} \deleted{\unsure{sans}} \unsure{sans pts fixes}}
% slanted
\unsure{sans} ramification est à la fois ouvert et fermé.

(ii) \deleted{\unsure{$X/G$ et $X$}} $X/G$ existe, et $X\to X/G$
et $X/G\mathrel{\unsure{-|\to}}Y$ sont \uline{étales}.

Si $s\in G$ \deleted{\unsure{qui fixe}} est dans le groupe
d'inertie d'un pt $x\in X$, il \deleted{laisse fixe} \inserted{opère \unsure{triv.|laisse} sur \unsure{les|un pt}} \unsure{la}
comp. \unsure{str.|strte} \unsure{connexe|connexes} de $x$ (\unsure{cf|sic|ii}), donc l'ens $X(s)$
des $x\in X$ tels que $s\in G_{i}(x)$ est une
réunion de composantes connexes, donc ouvert et
fermé. Comme \deleted{\unsure{$X(s)$}} l'ens des pts où
$G_i(x)\ne e$ est la réunion des $X(s)$ $(s\ne e)$,
(i) s'ensuit. D'autre part, $X$ \unsure{étant|etant}
\unsure{quasi-affine|qu. affine} sur $Y$, $X/G$ existe $[\,\unsure{\text{---}}\,]$
\deleted{\unsure{donc}} \unsure{on va prouver|on se borne au|en faisant} \unsure{ce qui était|cas fini|ce que} \unsure{.}

Soit $\uline{I}$ l'ens des comp. connexes de $X$.
$G$ opère sur $\uline{I}$, \deleted{\unsure{et est réuni au cas|et $X$ est réunion des}} \inserted{et $X$ est réunion des}
\unsure{des} \deleted{\unsure{$X(\alpha)$|$X_{(\alpha)}$}} $\alpha\in I/G$, obtenus en
\deleted{\unsure{par.|posant} $\unsure{\alpha}$} $X(\alpha)=\bigcup_{i\in\alpha}X_{i}$, stable sous $G$.
donc $X/G=\bigcup X(\alpha)/G$ on est
% archive sheet number: 156

\scanpage{158}
\uline{Catégorie galoisienne}

(i) Une catégorie $\unsure{\mathcal C|\mathscr C}$

(ii) Un foncteur $F:\mathcal C\to\text{Ensembles finis}$

Axiomes :

\margin{Poser plutôt en axiome l'\uline{existence} d'un $F$ ayant les \uuline{propriétés} : \deleted{\unsure{Soit|unicité}} $F$ permute aux sommes, aux produits, au \unsure{passage au quotient par un} groupe fini, et est tel qu'un \unsure{isomorphisme|injection} $F(X)\mathrel{\unsure{-|\to}} F(Y)$ \unsure{$X$ induit.|induit.|$X$ isom.} $\Rightarrow X\xrightarrow{\sim}Y$.}

(Gal 1) Dans $\mathcal C$, \deleted{\unsure{le|les}} sommes et \deleted{\unsure{le|les}} produits \deleted{de deux objets} \inserted{finis} existent
\inserted{[d'où l'existence de $\underline{0}$ et $\underline{1}$]}

\deleted{(Gal 3) $X\in\mathcal C$, $F(X)=\varnothing\Rightarrow X$ est une unité à gauche (soit $\underline{0}$)}

(Gal 4) $F$ permute à la somme et au produit

% margin note written sideways
\margin{Gal 3 et Gal 5 ensemble sont \unsure{impliqués} par : Si $f:X\to Y$ est tel que $F(X)\mathrel{\unsure{-|\to}}F(Y)$ soit \uline{injectif} \deleted{\unsure{($Y$ connexe)}} \inserted{($X$ \unsure{et|,} \deleted{\unsure{$Y$}} \unsure{connexe|connexes})} alors c'est un isomorphisme \inserted{de} $X$ sur une composante connexe \unsure{$Y_i$}.}

$X\in\mathcal C$ est dit \deleted{\uline{\unsure{irréductible|indivisible|indécomposable}}} \inserted{connexe} si $X=X'\amalg X''$ implique $X'=\underline{0}$
ou $X''=\underline{0}$

(Gal 2) $X\in\mathcal C$ \deleted{\unsure{$X\ne\underline{0}$}} implique $X\simeq\coprod_{i\in I}X_i$, \inserted{$X_i\ne\underline{0}$}, et tte décomposition
$X\simeq X'\amalg X''$ \unsure{impliques|implique} $X'\simeq\coprod_{i\in I'}X_i$, $X''=\coprod_{i\in I''}X_i$,
$I'\cap I''=\varnothing$, $I'\cup I''=I$. Les $X_i$, déterminés de
façon essentiellement unique [avec leur $X_i\to X$]
sont appelés les \uline{comp.} \deleted{\uline{\unsure{irréductibles|indivisibles|indécomposables}}} \inserted{connexes} de $X'$ [N.B. $X=\underline{0}\Leftrightarrow I=\varnothing$]

(\uuline{Gal 3}) Soit $f:X\to Y$ avec $Y$ \deleted{\unsure{irréductible|indivisible|indécomposable}} \inserted{connexe $\ne\underline{0}$}. Alors les
composantes \inserted{connexes} $X_i$ de $X$ telles que $f$ induise
un \uline{isom.} $f_i:X_i\xrightarrow{\sim}Y$ sont en correspondance
biunivoque avec les \uline{sections} $g_i:Y\to X$ \deleted{\unsure{$f$ de|de $f$}}
[hom. tels que $fg=\mathrm{id}_Y$], $g_i$ étant le composé
\deleted{\unsure{$Y\,X$}} $Y\xrightarrow{f_i^{-1}}X_i\xrightarrow{\text{inj. can}}X$ [en d'autres termes,
\deleted{toute} section se factorise] [N.B. il suffit
de le vérifier si $X=Z\times Y$, $f=p_2$]

(Gal 5) Sous les conditions de Gal 3, si \deleted{\unsure{$X$ est}} \inserted{$X,Y$ connexes et} $f:X\to Y$
induit $F(X)\xrightarrow{\sim}F(Y)$, c'est un isomorphisme
[implique $F(X)=\varnothing\to X=\underline{0}$] [Il suffit
de vérifier Gal 5 quand $X$ est une composante
connexe d'un produit $Z\times Y$, $f$ induit par $p_2$]

(\uuline{Gal 6}) Soit $X\in\mathcal C$, $X$ connexe $\ne\underline{0}$, et \uline{galoisien}, \deleted{\unsure{soit|i.e.}} i.e.
\deleted{\unsure{pour|Alors}} $G=\operatorname{Aut}(X)$ opère \uline{transitivement} dans $F(X)$
\deleted{\unsure{Tte composante connexe d'une puissance finie de} \illegible{1}}
Alors pour tout \unsure{ss-groupe} $H$ de $G$, $X/H$ dans $\mathcal C$ \uline{existe},
et $F(X/H)\simeq F(X)/H$
% archive sheet number: 157

\scanpage{159}
Catégorie $\unsure{\mathcal C|\mathscr C}$, sous-catégorie $\unsure{\mathscr P|\mathcal P}$\deleted{\unsure{,}} (pleine), morphismes ``non ramifiés''

\uline{1°} Dans $\mathcal C$, les sommes finies existent, ainsi que les produits fibrés.

\uline{4°} Dans $\mathcal C$, le \deleted{produit} \inserted{composé} de deux morphismes non ramifiés, \inserted{les morphismes identiques} ou l'image
inverse d'un morphisme non ramifié, sont non ramifiés, \deleted{\unsure{id. pour}}

\uline{5°} \deleted{\unsure{si}} $f:X\to Y$ est non ramifié $\Longleftrightarrow$ les $\unsure{f_i|f}\vert X_i$ ($X_i$ les comp. \unsure{connexes|conn.}) \deleted{\unsure{la somme|les sommes} $\coprod \unsure{X_i'}\to Y$} non ramifiés

\uline{2°} Dans $\mathcal C$, il y a des composantes connexes.

\uline{3°} Les \deleted{\unsure{pts|points}} $\xi\in\mathscr P$ sont connexes (donc $\operatorname{Hom}(\xi,X\amalg Y)=\operatorname{Hom}(\xi,X)\amalg\operatorname{Hom}(\xi,Y)$]

\uline{6°} Pour tout $f:X\to Y$ non ramifié, \inserted{$Y$ connexe,} et $\xi\in\mathscr P$, $\operatorname{Hom}(\xi,X)\mathrel{\unsure{-|\to}}\operatorname{Hom}(\xi,Y)$ est
\uline{surjectif}

\uline{7°} Si \deleted{\unsure{$X$ et}} $Y$ sont connexes, et la fibre de $\operatorname{Hom}(\xi,X)\to\operatorname{Hom}(\xi,Y)$ est
réduite à un élément, alors \deleted{\unsure{$X$}} $f$ est un \uline{isom}
(on suppose $f$ non ramifié, $\xi\in\mathscr P$)

\uline{8°} Si $f:X\to Y$ non ramifié, $Y$ connexe, alors \deleted{pour} toute
\deleted{\unsure{les}} section $g:Y\to X$ est un \uline{isom.} de $Y$ sur une
composante connexe de $X$.

\uline{9°} Soit \deleted{$G$ \inserted{\deleted{groupe fini}} opère dans $X$ ``sans poi} $f:X\to Y$ \uline{ramifié}, et
$X$ connexe, \deleted{\unsure{$\xi\in$}} $e\in\operatorname{Hom}(\xi,Y)=\unsure{:|\,}Y^\xi$, supposons $G=\operatorname{Aut}_Y(X)$
opère transitiv\unsure{ement|emt} dans $(f^\xi)^{-1}(e)$. Alors pour tout
sous-groupe $H$ de $G$, $X/H$ existe \inserted{dans $\mathcal C$} et est un
revêtement non ramifié de $Y$. \deleted{et $X$ est un \unsure{rev. princ. pour} $H$,}
i.e. deux élts de $\operatorname{Hom}(\xi,X)$ \unsure{au dessus de} $e$
ayant \unsure{la même image} par \deleted{\unsure{$X$}} $\mathrel{-} X/H$ sont \unsure{congrus|conjugués} mod $H$.

Sous ces conditions, pour $x\in\operatorname{Hom}(\xi,X)$ on définit
$\pi_1(X,x)$. Il ne change pas \deleted{\unsure{si}} en remplaçant $\xi$ par $\xi'$ tel que
$\xi'\mathrel{-}\xi$ soit donné. C'est un foncteur \uline{covariant} en $(X,x)$

\diagram{\begin{tikzcd}[ampersand replacement=\&]
\xi \arrow[r,dash,"x"] \arrow[d] \& X \arrow[d]\\
\xi' \arrow[r,dash,"x'"] \& X'
\end{tikzcd}}
définissant $\pi_1(X,x)\mathrel{-}\pi_1(X',x')$.

\uline{Etc.}
% archive sheet number: 158

\scanpage{160}
Groupe fondamental dans les catégories.

$\unsure{\mathcal C|\mathscr C}$ une catégorie, avec une famille $(D)$ de morphismes
de $\mathcal C$, satisfaisant les conditions usuelles

(i) contient identité et stable par composition (ii) stable par
changement de base

On suppose de plus

(iii) \deleted{\unsure{les}} \inserted{$D$} sont des morphismes de descente [donc des
\unsure{épim.} stricts \inserted{universels} si \inserted{\unsure{on suppose que}} \uline{les produits et produits}
\uline{fibrés existent}]. Enfin on suppose

% slanted
\margin{On peut \unsure{axer|varier|aussi} la situation, en introduisant de plus \deleted{\unsure{la}} une classe de morphismes $D'$ [\deleted{de} morphismes pouvant être \unsure{``localement''} \unsure{fibrés}.}

(iv) On peut ``descendre'' un $g\in(D)$ par un $f\in(D)$
[mais on ne suppose pas qu'on peut descendre
\uline{tout} morphisme par un $f\in(D)$]. \uline{Dans la
suite, $D$ est supposé donné une fois pour toutes.}
\uline{On suppose que dans $\mathcal C$ les sommes finies} [\unsure{resp.|évent.} quelconques] existent\unsure{]|.}

\margin{(v) La projection $X\times I\to X$ est $\in(D)$.}

\uline{Exemple} Catégorie \inserted{et \unsure{première?} condition} $\unsure{\mathscr E|\mathcal E|\mathcal C}$ des préschémas. \deleted{La catégorie}
Classe $(D)$ des morphismes \uline{plats} et \uline{affines}.
[ou \uline{plats} et quasi-projectifs].

On va développer une théorie particulière des
fibrés principaux [vérifiant les conditions générales
\deleted{\unsure{Groupes de base|Groupes de faisceaux}}
d'une théorie des fibrés principaux, cf. \unsure{---}].

\uline{Groupes structuraux} : $X\times G$ $G$ un groupe ensembliste
fini [\unsure{resp.} quelconque]

\uline{Fibrés principaux \unsure{locaux} sur $X$}, de groupe $G$ : les
fibrés \uline{\unsure{fonctoriels|formels|fonctls|formés}} ainsi, avec hom $P\mathrel{\unsure{-|\to}}X$ $\in(D)$
% archive sheet number: 159

\scanpage{161}
\uline{\unsure{Espace|Espaces} à opérateurs admissibles} : \unsure{Le|Les} \unsure{morphisme|morphismes} $E\to X$ est $\in D'$.

\uline{Hom} \uline{admissibles} de \deleted{groupes} \uline{faisceaux structuraux} : déduits
\uline{d'un hom de groupes usuels}

\deleted{Hom} \uline{Inclusions admissibles} : \deleted{\unsure{les|hom}} monomorphismes
de groupes ordinaires.
% archive sheet number: 160

\scanpage{162}
Soit $f:X\mathrel{\unsure{-|\to}}Y$ un morphisme \inserted{propre} dominant, $Y$ \inserted{loc. noeth.} \uline{normal} \uline{connexe}.
Alors $\pi_1(X)\to\pi_1(Y)$ a un conoyau fini (regardant
le pt générique $y$ de $Y$, et notant que $\pi_1(y)\to\pi_1(Y)$ est
surjectif, on est ramené au cas où $Y=\operatorname{Spec}(K)$.
On voit de \unsure{$=$|même} qu'on peut supposer $X$ séparable sur $k$
et \unsure{abst|absolt} connexe [remplacer $K$ par une extension \inserted{séparable}
finie convenable $K'$], et diviser \deleted{\unsure{$O$}} $X'$ par ses élts
nilpotents, enfin prendre une composante connexe
de $X'$]. Mais alors $\pi_1(X)\to\pi_1(Y)$ est \uline{surjectif}..

\noindent\rule{2em}{0.25pt}

N.B. Si $Y$ n'est pas \uline{normal}, le th est faux,
\unsure{$=$|même|si} si p.ex $X$ est le normalisé de $Y$. Ex. $Y$ courbe
\unsure{rat.} sur $k$ alg. clos, avec pt double ordinaire,
$X'$ sa normalisée. On a $\pi_1(X)=0$, $\pi_1(Y)=\widehat{\mathbb Z}$
privé de \deleted{\unsure{sa}} $p$-composante.

\noindent\rule{2em}{0.25pt}

On voit facilement que dans le dernier
\unsure{exemple|exemples}, on ne peut former $\underline{\operatorname{Hom}}_Y(X,Y\amalg Y)$
[car sur un \inserted{\unsure{certain}} pt $y_0$ de $Y$, il y a deux pts
géométriques dans l'image inverse, alors
qu'il n'y en a qu'un au voisinage].

\uline{Supposons cependant} $f:\unsure{X|x}\to Y$ \uline{universellement}
\uline{ouvert} [p. ex. $Y$ \unsure{normale|unibranche}], je crois qu'on
% archive sheet number: 161

\scanpage{163}
peut alors \deleted{\unsure{des fai|des fo|du fai}} construire les $f_*(X_1)$ pour
$X_1$ étale sur $X$. \unsure{Partons|Partir|P.-ê.} plutôt\unsure{)|:}

\diagram{\fbox{\begin{tikzcd}[ampersand replacement=\&,column sep=small,row sep=tiny]
\& X \arrow[d]\\
S \& T \arrow[l,very thick]
\end{tikzcd}}}
% archive sheet number: 162

\scanpage{164}
\deleted{\unsure{pro Groupe|proGroupe|Le Groupe}}

\uline{Progroupe fondamental}
\uline{pour les revêtements}
\uline{infinis} \uline{principaux}
% archive sheet number: 163

\scanpage{165}
Le $\widetilde{\pi}_1(S,\xi)$, $S$ \unsure{préschéma|préschém.} connexe, $\xi$ pt géom.

Pour tout groupe $\pi$, soit
\[
\unsure{\Pi|\pi}^{1}(S,\unsure{\xi|\vert};\pi)
\]
l'ens des classes, à isom près, de \deleted{\unsure{rev.|esp.}} \inserted{fibrés} principaux
homogènes sur $S$, de groupe structural $\pi$,
ponctués au dessus de $\xi$.

(i) $\pi\leadsto\Pi^1(S,\xi;\pi)$ est exact à gauche.

En effet \quad a) Commute aux produits finis
c'est trivial

b) Transforme monos \inserted{$\pi\mathrel{\unsure{-|\to}}\pi'$} en \unsure{mono.|mon.}

En effet, cela provient que \deleted{\unsure{pour tt}} \inserted{si}
\deleted{isom} $P'_1\mathrel{\overset{u'}{\unsure{-|\to}}}P'_2$ \deleted{\unsure{des|de}} \inserted{deux} $\pi'$ fibrés \deleted{\unsure{est}} isom,
il existe un et un seul tel isom
qui applique $\xi'_1$ en $\xi'_2$. Si alors
$P'_1$ et $P'_2$ sont \uline{associés} à $P_1,P_2$ d'où
\[
\left\{\begin{aligned}P_i&\subset P'_i\\ \xi_i&=\xi'_i\end{aligned}\right.,
\]
\deleted{\unsure{donc}} je dis que $\unsure{u|u'}$ \deleted{applique}
\deleted{\unsure{$P_i$}} induit un isom $u:P_1\mathrel{-}P_2$. En
effet, \deleted{$P_2$ est ouvert et fermé dans $P'_2$ et \unsure{contient} $\xi'_2$,}
$\unsure{u^{-1}|{u'}^{-1}}(P_2)$ est \uline{ouvert} et \uline{fermé} dans
$P'_1$, et contient $\xi$, enfin il est \unsure{inv. pour|invariant pour|sous-groupe}
% archive sheet number: 164

\scanpage{166}
\deleted{\unsure{$\pi$. Je}} $P_i$ se déduit de $P'_i$ par une
\uline{section} de $P'_i/\pi$ passant par $\overline{\xi'}$. $S$ \unsure{étant|etant}
connexe, il ne peut y en avoir qu'une
seule telle section.

c) Si $\pi\mathrel{\unsure{-|\to}}\pi'\rightrightarrows\pi''$ est exact,
alors \deleted{\unsure{$H(\,$}} $\unsure{\Phi|\varphi}(\pi)\mathrel{-}\Phi(\pi')\rightrightarrows\Phi(\pi'')$ aussi.

Il \deleted{\unsure{reste|veux}} à prouver que si $P'$ de groupe $\pi'$,
pt marqué $\xi'$, est tel que \deleted{\unsure{les}} l'on
ait un isom $P'\times_{\pi'}(\pi'',\varphi_1)\overset{u}{\xrightarrow{\sim}}P'\mathbin{\unsure{\times|*}}_{\pi'}(\pi'',\varphi_2)$,
alors on peut restreindre le groupe structural.

Or soit $Q$ le sous-schéma de $\unsure{P|P'}$
des coïncidences de $ui_1$ et $i_2$.

% diagram in the left margin, above the triangle
\diagram{\begin{tikzcd}[ampersand replacement=\&,column sep=small]
\pi \arrow[r] \& \pi' \arrow[r,shift left=.45ex] \arrow[r,shift right=.45ex] \& \pi''
\end{tikzcd}}
\diagram{\begin{tikzcd}[ampersand replacement=\&,column sep=small,row sep=small]
P' \arrow[r,"i_1"] \arrow[dr,"i_2"'] \& P''_1 \arrow[d,"u"]\\
\& P''_2
\end{tikzcd}}

Il est ouvert et fermé, contient
$\xi'$, est stable par $\pi$, donc
correspond à un sous-schéma de
$P'/\pi$ à la fois ouvert et fermé $S'$
[descente \inserted{essentiellement} \unsure{fidèlt|fidèlement} plat et \unsure{quasi-compacte|quasi-compact} \unsure{---|...}]
De plus, on voit que les \deleted{\unsure{degrés}} fibres géom
de $S'\mathrel{-}S$ sont partout un pt ; donc
c'est un isom $S'\to S$ \unsure{..}, on a gagné\unsure{...}
% archive sheet number: 165

\scanpage{167}
\inserted{Supposons \deleted{$S$} $S$ loc. noeth \uline{ou} $\pi$ artinien [condition des chaînes descendantes]}

(ii) Tout élément $c\in\unsure{\Phi|\varphi}(\pi)$ est \unsure{majoré|minoré} par
un élément \unsure{minimal}, $c'\in\Phi(\pi')$.

% slanted
\margin{\unsure{par un} \unsure{ss-gr.} \unsure{et} \illegible{1} \uline{\unsure{est-il|minimal}} \unsure{NB|2e}}

En effet, si $c$ correspond à $P^{\unsure{\pi}},y$,
supposons que l'on puisse \unsure{réduire|relever} à $P',y'$ \inserted{\unsure{$P'$ (groupe $\pi'\subset\pi$)}} \deleted{\unsure{considérant les images connexes de} \unsure{$y$}.}
et $(P'',y'')$ [groupe $\pi''\subset\pi$].

\margin{$\pi/{\pi'\cap\pi''}\to\pi/\pi'\times\pi/\pi''$}
Considérons $P/{\pi'\cap\pi''}\to P/\pi'\times P/\pi''$
c'est une imm. ouverte et fermée, d'autre
part on a une section du $2^{\mathrm{o}}$ \unsure{membre}
qui \deleted{est} au pt $(y',y'')$ est dans l'image
du premier, elle y est donc partout.

Donc on peut restreindre le groupe structural
à $\pi'\cap\pi''$. Donc les groupes auxquels
on peut restreindre forment un ens
filtrant décroissant, $\pi_i,\ \unsure{\ldots}$, donc $P_i\subset P$
\inserted{forment un syst. projectif.}

Si \deleted{\unsure{les morphismes}} $\pi$ est artinien, on a gagné : le \unsure{systême|système} des $\pi_i$ est
stationnaire. Si $S$ est loc. noeth, on
\deleted{\unsure{les morphismes} de transition $P_j\mathrel{-}P_i$ sont
affines. Donc on peut former le
préschéma $\varprojlim_i P_i$, qui est donc affine sur
$P$, \deleted{\unsure{limite,} \unsure{il est}} \inserted{je dis que l'un des \unsure{objets|sous-groupes}} est
ouvert et fermé, et est un fibré}
% archive sheet number: 166

\scanpage{168}
\deleted{principal homogène sous $\pi'=\bigcap\pi_i$. On voit
ici que la question est locale sur $S$
[par \unsure{prop. de trans.|recollement|propagation} le long d'ouverts affines],
donc on peut supposer \uline{$S$ affine}. Mais alors
il existe $S_{\unsure{0|o}}^{(0)}$ de type fini sur $\mathbb Z$, \inserted{connexe,} et
$(P^{(0)},\xi^{(0)})$ principal homogène \inserted{ponctué} sur $\pi_{\unsure{S^{(0)}|S(0)}}$, qui
donne $P$ par extension de la base [se rappeler
que $P$ est trivialisé par $S'\mathrel{\unsure{-|\to}}S$ quasi-fini
étale !]. On peut supposer $\unsure{S|S'}$ \uline{connexe}.}

\diagram{\begin{tikzcd}[ampersand replacement=\&,column sep=small,row sep=small]
P^{0} \arrow[d,dash] \& P \arrow[l] \arrow[d]\\
S^{0} \& S \arrow[l]
\end{tikzcd}}

\deleted{Alors les}

\unsure{l'en tire|l'on tire|s'en tire} à \unsure{meilleur compte|meilleurs comptes|meilleur résultat} : on
regarde la composante connexe $P^{0}$ de $\xi$ dans $P$
[\deleted{existe et} est \uline{\unsure{ouverte|ouvert} et fermée}, car $P$ est
loc noeth donc loc connexe]. Je dis que
$P^0\to S$ est \uline{surjectif}. Comme l'image est
ouverte, pour prouver qu'elle est fermée,
on est ramené au cas d'un anneau
de valuation discrète complet \inserted{à corps résiduel alg clos} : à prouver que
si $P$ est un fibré principal \unsure{hom|hom.} sur $S$
% archive sheet number: 167

\scanpage{169}
\uline{Application} Un schéma étale \inserted{$T$ \uline{sur $S$}} est dit à \uline{groupe}
\uline{structural} $\pi$, \uline{fibre $F$}, si on a
\deleted{\unsure{si on se donne un} sous-schéma de $\underline{\operatorname{Aut}}_S(T)$, qui est}
donné \unsure{une opération fidèle|des opérations fidèles} de $\pi$ sur $F$,
et un sous-fibré \inserted{$\pi_S$} principal de
$\underline{\operatorname{Isom}}_S(E_S,T)$. Si $S$ est connexe et
muni du pt géométrique $\xi$, il revient
au $=$ de se donner un $T$ avec
une telle structure, ou de se
donner des opérations \deleted{continues}
de $\widetilde{\pi}_1(S,\xi)$ \deleted{\unsure{à travers}} sur $F$ à
travers $\widetilde{\pi}_1(S,\xi)\unsure{:|\,}\to\deleted{\unsure{G}}\pi$.

\uline{Exemple} Schémas en groupes comm.
\unsure{constants} tordus sur $S$, \deleted{à fibres \unsure{gé.}} \inserted{de}
génération finie sur $S$, correspondent
\deleted{\unsure{aux|au}} $(\widetilde{\pi}_1(S,\xi))$-Modules admissibles
\deleted{de \unsure{présentat|présentation}} qui sont des $\mathbb Z$-modules
de type fini \unsure{---}
% archive sheet number: 168

\scanpage{170}
de groupe un groupe constant $\pi_S$, alors
\uline{tte composante connexe} $P^0$ se surjecte sur $S$.
Or c'est trivial, car $P$ est \uline{trivial} !

Ceci étant, comme $P^{(0)}\subset\deleted{\unsure{P|\bigcap}}P^{(i)}$ pour tt $i$,
on voit que $\bigcap_i P^{(i)}$ \deleted{\unsure{contient}} un pt
de chaque fibre. \deleted{\unsure{Comme} les inclusions \unsure{de}}
\deleted{\unsure{Il s'ensuit que}}
\deleted{\unsure{fibré des|finie}} \inserted{chaque} $P^{(i)}$ est \unsure{une réunion|un réunion} de
composantes connexes de $P$, il en est de
\unsure{$=$|même} de l'intersection, \deleted{qui} $P'$, qui contient
$P^0$. Elle est stable \unsure{sous|sur} $\pi'$, et on \unsure{conclut|constate}
aussitôt (en regardant sur chaque fibre
géométrique, et utilisant que $P'_{\unsure{\bar s|s}}\ne\varnothing$]
que $P'$ est fibré principal homogène \unsure{sur|sous} $\pi'$.
cqfd.

\uline{Conséquences} Si $\pi$ varie dans la \inserted{\unsure{une|une sous-}} catégorie
des groupes artiniens, \deleted{\unsure{ou si $S$ est}}
stable par produits finis et sous-groupes,
ou si $\pi$ varie dans n'importe quelle
catégorie stable par ces opérations, mais $S$
\unsure{étant|etant} loc noeth --- alors le foncteur
$\pi\leadsto\unsure{\Pi|\pi}^{1}(S,\xi;\pi)$ est strictement proreprésentable\unsure{$^{*}$|/}

% margin note written sideways
\margin{\unsure{$^{*}$} par un progroupe $\widetilde{\pi}_1(S,\xi)$ [qui ne peut plus être interprété comme un simple groupe topologique \unsure{---}]}
% archive sheet number: 169

\scanpage{171}
\uline{Groupe fondamental} : (Revêtements infinis)

$S$ préschéma. Pour tout groupe discret $G$, on définit
\[
H^1_{\unsure{\mathcal C|\mathscr C}}(S,G)=\varinjlim_{\substack{T/S\in\mathcal C\\ \deleted{\unsure{T\ \text{fid. plat de type fini sur }S}}\\ \deleted{\unsure{\text{de type qct sur }S|\text{de type qct vers }S}}}}H^1(T/S,G)
\]
où
\[
H^1(T/S,G)=H^1(K_{T/S},G)=H^1(\pi_0(K_{T/S}),G)
\]
et où $\mathcal C$ est une catégorie de $S$-préschémas.

% slanted
\margin{élucider ici les questions d'effectivité des données de descente\unsure{.....} \unsure{Se ramener|Ça ramène} à la descente des morphismes étales de type fini\unsure{,....}}

Les cas les plus importants sont les suivants

\diagram{\(\begin{array}{l}
\left\{\begin{array}{l}
\text{(i)}\quad \mathcal C_1\qquad S\text{-préschémas fid plats quasi-compacts, }\deleted{\unsure{\text{localemt}|\text{localement}}}.\\
\left.\begin{array}{lll}
\text{(ii)}&\mathcal C_2&\rule{4em}{0.3pt}\quad\text{de type fini}\\
\text{(iii)}&\mathcal C_3&\rule{4em}{0.3pt}\quad\text{quasi-finis}
\end{array}\right\}
\end{array}\right.\\[3pt]
\left.\begin{array}{lll}
\text{(iv)}&\mathcal C_4&\rule{4em}{0.3pt}\quad\text{finis par morceaux}\\
\text{(v)}&\mathcal C_5&\rule{4em}{0.3pt}\deleted{\text{étales}}\quad\text{finis étales par morceaux}
\end{array}\right\}
\end{array}\)}

On a donc $H^1_{\mathcal C_i}(\deleted{\unsure{T,|P}}S,G)\hookrightarrow H^1_{\mathcal C_j}(S,G)$ si $i\geq j$
application injective canonique. \uline{Si $G$ est fini},
alors tous les hom sont bijectifs, et $H^1(S,G)$
est l'ensemble des classes (à isom près) des
revêtements étales galoisiens de groupe $G$, donc $\simeq$
$\operatorname{Hom}(\pi_1(S,a),G)/\operatorname{int}(G)$. Notre intérêt ici est pour
le cas $G$ infini. Notons en passant
\[
H^1_{\mathcal C_3}\xrightarrow{\sim}H^1_{\mathcal C_2}
\]
% archive sheet number: 170

\scanpage{172}
et si $S$ est le spectre d'un anneau de Dedekind
\[
H^1_{\unsure{\mathcal C|\mathscr C|C}_2}\xrightarrow{\sim}H^1_{\mathcal C_1}
\]
[d'où dans ce cas $H^1_{\mathcal C_3}\xrightarrow{\sim}H^1_{\mathcal C_2}\simeq H^1_{\mathcal C_1}$]. Si $S$ est le spectre d'un anneau \deleted{local complet}\inserted{de valuation} discrète à corps résiduel alg. clos, \inserted{$V$} on aura de $\overline m$
\[
H^1_{\mathcal C_i}\unsure{b|6}\simeq H^1_{\deleted{\unsure{\mathcal C_i|4}}}(\widehat V/V,-)\qquad\text{pour }i=1,2,3
\]
Si $S$ est le spectre d'un anneau local \uline{complet} $V_{\unsure{\phi|0}}$, alors
\[
\deleted{H^1_{\unsure{\mathcal C_4|4}}(\unsure{V|k},\unsure{-|\,})\simeq}
\]
\[
H^1_{\mathcal C_5}\simeq H^1_{\mathcal C_4}\simeq H^1_{\mathcal C_3}\quad(\simeq H^1_{\mathcal C_2})
\]
et on aura de plus
\[
H^1_{\mathcal C_i}(V,-)\simeq H^1_{\mathcal C_i}(\underset{\substack{\downarrow\\\text{corps}\\\text{résiduel}}}{k},-)\qquad(i=2,3,4,5)
\]
% archive sheet number: 171
